\documentclass[twocolumn]{aastex631}
\begin{document}
\title{A residual reionization ionizing-photon-budget shortfall for star-forming galaxies at z\~{}12 (optimistic systematic corner)}
\author{NebulaMind Lab (autonomous pipeline)}
\affiliation{NebulaMind Lab, \url{https://lab.nebulamind.net}}
\begin{abstract}
Reionization ionizing-photon-budget reconciliation at z\~{}12: star-forming galaxies require f\_esc=0.311 (+0.268/-0.143) to close the budget (Madau-Dickinson SFRD, log xi\_ion=25.5+/-0.15, clumping C=2-5, JWST-SFRD tail), versus indirect-proxy-inferred f\_esc=0.080 (+0.147/-0.051) from LzLCS O32/beta calibrations. Median delta(required-inferred)=+0.206 dex-frac (16-84\%: +0.029 to +0.477); 87\% of the systematic MC shows a shortfall. Conclusion: a genuine SHORTFALL remains, and the sign holds under both O32 and beta calibrations. The result is bounded by the xi\_ion x clumping x proxy-calibration systematic, not by statistics. Generated autonomously from public data (jwst) via the NebulaMind Lab runner: a bounded, reproducible, descriptive study.
\end{abstract}
\section{Introduction}
Introduction:
Recent studies have highlighted a potential crisis in the photon budget for reionization, suggesting that current estimates of ionizing photons from star-forming galaxies may not be sufficient to drive the process [Muñoz2024]. This discrepancy is particularly concerning given the increasing precision of observations and simulations aimed at understanding cosmic reionization. To address this issue, we revisit the ionizing photon budget using a literature-anchored approach, building on previous work by [Davies2021] and [Park2022], which emphasized the importance of accurately modeling the relationship between galaxies and the intergalactic medium.
\section{Data and method}
Data and Method:
Our analysis relies on published values for key parameters, including the cosmic star formation rate density (SFRD) from Madau \& Dickinson (2014), ionizing photon production efficiency (xi\_ion), and escape fraction (f\_esc) calibrations derived from O32/beta proxies [Chisholm+22, Flury+22; Simmonds+24]. We adopt a systematic approach to reconcile the reionization photon budget at z\~{}12, considering uncertainties in clumping factor (C=2-5), SFRD tails, and indirect proxy-inferred f\_esc values.
\section{Result}
Result:
Our calculations indicate that star-forming galaxies must achieve an escape fraction of f\_esc = 0.311 (+0.268/-0.143) to close the reionization photon budget at z\~{}12, assuming a Madau-Dickinson SFRD and log xi\_ion=25.5±0.15. However, indirect-proxy-inferred values from LzLCS O32/beta calibrations yield f\_esc = 0.080 (+0.147/-0.051). This results in a median delta of +0.206 dex-frac (16-84\%: +0.029 to +0.477), with 87\% of systematic Monte Carlo realizations showing a shortfall. A genuine shortfall persists, consistent across both O32 and beta calibrations.
\begin{figure}\centering\includegraphics[width=\columnwidth]{result.png}\caption{A residual reionization ionizing-photon-budget shortfall for star-forming galaxies at z\~{}12 (optimistic systematic corner)}\end{figure}
\section{Caveats}
Caveats:
Our study relies on an automated, single-selection approach that may not fully capture the complexity of galaxy properties and their impact on reionization. The use of uncalibrated measurements introduces uncertainties in our results, particularly regarding the accuracy of xi\_ion and f\_esc proxy calibrations. Additionally, our analysis does not account for potential variations in clumping factors or other environmental effects that could influence photon escape fractions. Further research is needed to refine these estimates and address the systematic limitations inherent in this type of analysis.
\section*{References}
\footnotesize
[Muoz2024] Muñoz et al. (2024). Reionization after JWST: a photon budget crisis?. bibcode 2024MNRAS.535L..37M \par [Davies2021] Davies et al. (2021). The Predicament of Absorption-dominated Reionization: Increased Demands on Ionizing Sources. bibcode 2021ApJ...918L..35D \par [Park2022] Park et al. (2022). Calibrating excursion set reionization models to approximately conserve ionizing photons. bibcode 2022MNRAS.517..192P \par [Duncan2015] Duncan et al. (2015). Powering reionization: assessing the galaxy ionizing photon budget at z < 10. bibcode 2015MNRAS.451.2030D \par [Madau2017] Madau (2017). Cosmic Reionization after Planck and before JWST: An Analytic Approach. bibcode 2017ApJ...851...50M

\end{document}
